
\documentclass{article}
\usepackage{colortbl}
\usepackage{makecell}
\usepackage{multirow}
\usepackage{supertabular}

\begin{document}

\newcounter{utterance}

\twocolumn

{ \footnotesize  \setcounter{utterance}{1}
\setlength{\tabcolsep}{0pt}
\begin{supertabular}{c@{$\;$}|p{.15\linewidth}@{}p{.15\linewidth}p{.15\linewidth}p{.15\linewidth}p{.15\linewidth}p{.15\linewidth}}

    \# & $\;$A & \multicolumn{4}{c}{Game Master} & $\;\:$B\\
    \hline 

    \theutterance \stepcounter{utterance}  

    & & \multicolumn{4}{p{0.6\linewidth}}{\cellcolor[rgb]{0.9,0.9,0.9}{%
	\makecell[{{p{\linewidth}}}]{% 
	  \tt {\tiny [A$\langle$GM]}  
	 Du nimmst an einem kollaborativen Verhandlungspiel Teil.\\ \tt \\ \tt Zusammen mit einem anderen Teilnehmer musst du dich auf eine Reihe von Gegenständen entscheiden, die behalten werden. Jeder von euch hat eine persönliche Verteilung über die Wichtigkeit der einzelnen Gegenstände. Jeder von euch hat eine eigene Meinung darüber, wie wichtig jeder einzelne Gegenstand ist (Gegenstandswichtigkeit). Du kennst die Wichtigkeitsverteilung des anderen Spielers nicht. Zusätzlich siehst du, wie viel Aufwand jeder Gegenstand verursacht.  \\ \tt Ihr dürft euch nur auf eine Reihe von Gegenständen einigen, wenn der Gesamtaufwand der ausgewählten Gegenstände den Maximalaufwand nicht überschreitet:\\ \tt \\ \tt Maximalaufwand = 3624\\ \tt \\ \tt Hier sind die einzelnen Aufwände der Gegenstände:\\ \tt \\ \tt Aufwand der Gegenstände = {"C33": 210, "A89": 291, "A48": 831, "A04": 111, "A41": 926, "C64": 827, "C49": 25, "A61": 121, "C04": 583, "B53": 766, "B85": 14, "B61": 559, "A45": 304, "A03": 989, "B21": 691}\\ \tt \\ \tt Hier ist deine persönliche Verteilung der Wichtigkeit der einzelnen Gegenstände:\\ \tt \\ \tt Werte der Gegenstandswichtigkeit = {"C33": 138, "A89": 583, "A48": 868, "A04": 822, "A41": 783, "C64": 65, "C49": 262, "A61": 121, "C04": 508, "B53": 780, "B85": 461, "B61": 484, "A45": 668, "A03": 389, "B21": 808}\\ \tt \\ \tt Ziel:\\ \tt \\ \tt Dein Ziel ist es, eine Reihe von Gegenständen auszuhandeln, die dir möglichst viel bringt (d. h. Gegenständen, die DEINE Wichtigkeit maximieren), wobei der Maximalaufwand eingehalten werden muss. Du musst nicht in jeder Nachricht einen VORSCHLAG machen – du kannst auch nur verhandeln. Alle Taktiken sind erlaubt!\\ \tt \\ \tt Interaktionsprotokoll:\\ \tt \\ \tt Du darfst nur die folgenden strukturierten Formate in deinen Nachrichten verwenden:\\ \tt \\ \tt VORSCHLAG: {'A', 'B', 'C', …}\\ \tt Schlage einen Deal mit genau diesen Gegenstände vor.\\ \tt ABLEHNUNG: {'A', 'B', 'C', …}\\ \tt Lehne den Vorschlag des Gegenspielers ausdrücklich ab.\\ \tt ARGUMENT: {'...'}\\ \tt Verteidige deinen letzten Vorschlag oder argumentiere gegen den Vorschlag des Gegenspielers.\\ \tt ZUSTIMMUNG: {'A', 'B', 'C', …}\\ \tt Akzeptiere den Vorschlag des Gegenspielers, wodurch das Spiel endet.\\ \tt STRATEGISCHE ÜBERLEGUNGEN: {'...'}\\ \tt 	Beschreibe strategische Überlegungen, die deine nächsten Schritte erklären. Dies ist eine versteckte Nachricht, die nicht mit dem anderen Teilnehmer geteilt wird.\\ \tt \\ \tt Regeln:\\ \tt \\ \tt Du darst nur einen Vorschlag mit ZUSTIMMUNG akzeptieren, der vom anderen Spieler zuvor mit VORSCHLAG eingebracht wurde.\\ \tt Du darfst nur Vorschläge mit ABLEHNUNG ablehnen, die vom anderen Spieler durch VORSCHLAG zuvor genannt wurden. \\ \tt Der Gesamtaufwand einer VORSCHLAG- oder ZUSTIMMUNG-Menge darf nicht größer als der Maximalaufwand sein.  \\ \tt Offenbare deine versteckte Wichtigkeitsverteilung nicht.\\ \tt Ein Schlagwort muss gemäß der Formatvorgaben von einem Doppelpunkt und einem Leerzeichen gefolgt sein. Das Argument muss eine Python-Menge mit 0 oder mehr Strings sein.  \\ \tt Die Form muss also SCHLAGWORT: {...} sein.\\ \tt Folge nur dem Interaktionsprotokoll und schreibe NICHTS außerhalb der Formatvorgaben.\\ \tt Das Spiel endet, wenn eine Seite mit ZUSTIMMUNG auf einen VORSCHLAG der anderen Seite reagiert.  \\ \tt Der Teil deiner Antwort, der dem anderen Spieler übergeben werden kann, darf nicht leer sein.  \\ \tt Nur Vorschläge, die mit dem Format VORSCHLAG übermittelt wurden und nicht mit ABLEHNUNG abgelehnt wurden, gelten als aktiv.  \\ \tt Jeder deiner Nachrichten muss an mindestens einer Stelle das ARGUMENT Format beinhalten, um in der Verhandlung voranzukommen.\\ \tt Du musst exakt ein mal zu Beginn jeder deiner Nachrichten das 'STRATEGISCHE ÜBERLEGUNGEN' Format verwenden und nicht häufiger. Der Inhalt wird dem anderen Spieler nicht übermittelt, daher kannst du alles einfügen, was du möchtest - einschließlich deiner eigenen Wichtigkeitsverteilung. Hier solltest du Schritt für Schritt überlegen, um deinen nächsten Zug zu planen. 
	  } 
	   } 
	   } 
	 & \\ 
 

    \theutterance \stepcounter{utterance}  

    & \multicolumn{4}{p{0.6\linewidth}}{\cellcolor[rgb]{0.8,1,0.9}{%
	\makecell[{{p{\linewidth}}}]{% 
	  \tt {\tiny [A$\rangle$GM]}  
	 STRATEGISCHE ÜBERLEGUNGEN: {'Ich muss meine Wichtigkeit maximieren, während der Gesamtaufwand <= 3624 bleibt. Meine wichtigsten Gegenstände sind A48 (868), B21 (808), A04 (822), A41 (783), B53 (780) und A89 (583) mit hohen Werten. Aber ich muss auch Aufwände berücksichtigen: A48 (831), A41 (926), B53 (766), B21 (691), A89 (291) und A04 (111) sind teuer. Ich sollte versuchen, hochwertige, niedrigaufwändige Gegenstände wie A04 (822 Wert, 111 Aufwand) und B85 (461 Wert, 14 Aufwand) einzubeziehen, während teure, niedrigwertige Gegenstände wie A03 (389 Wert, 989 Aufwand) vermieden werden. Ich kenne die Präferenzen des anderen nicht, also sollte ich einen ausgeglichenen Vorschlag machen, der mir nützt, aber auch für den anderen akzeptabel sein könnte.'}\\ \tt \\ \tt VORSCHLAG: {'A04', 'A89', 'A48', 'A41', 'B53', 'B21', 'B85', 'C04'}  \\ \tt ARGUMENT: {'Dieser Vorschlag enthält wertvolle Gegenstände mit einem Gesamtaufwand von 3624, was genau dem Maximum entspricht. Er bietet eine gute Mischung aus verschiedenen Kategorien und sollte für beide Seiten vorteilhaft sein.'} 
	  } 
	   } 
	   } 
	 & & \\ 
 

    \theutterance \stepcounter{utterance}  

    & & \multicolumn{4}{p{0.6\linewidth}}{\cellcolor[rgb]{0.9,0.9,0.9}{%
	\makecell[{{p{\linewidth}}}]{% 
	  \tt {\tiny [GM$\rangle$B]}  
	 Du nimmst an einem kollaborativen Verhandlungspiel Teil.\\ \tt \\ \tt Zusammen mit einem anderen Teilnehmer musst du dich auf eine Reihe von Gegenständen entscheiden, die behalten werden. Jeder von euch hat eine persönliche Verteilung über die Wichtigkeit der einzelnen Gegenstände. Jeder von euch hat eine eigene Meinung darüber, wie wichtig jeder einzelne Gegenstand ist (Gegenstandswichtigkeit). Du kennst die Wichtigkeitsverteilung des anderen Spielers nicht. Zusätzlich siehst du, wie viel Aufwand jeder Gegenstand verursacht.  \\ \tt Ihr dürft euch nur auf eine Reihe von Gegenständen einigen, wenn der Gesamtaufwand der ausgewählten Gegenstände den Maximalaufwand nicht überschreitet:\\ \tt \\ \tt Maximalaufwand = 3624\\ \tt \\ \tt Hier sind die einzelnen Aufwände der Gegenstände:\\ \tt \\ \tt Aufwand der Gegenstände = {"C33": 210, "A89": 291, "A48": 831, "A04": 111, "A41": 926, "C64": 827, "C49": 25, "A61": 121, "C04": 583, "B53": 766, "B85": 14, "B61": 559, "A45": 304, "A03": 989, "B21": 691}\\ \tt \\ \tt Hier ist deine persönliche Verteilung der Wichtigkeit der einzelnen Gegenstände:\\ \tt \\ \tt Werte der Gegenstandswichtigkeit = {"A48": 65, "A04": 121, "B21": 138, "A41": 262, "B53": 389, "A45": 461, "A89": 484, "C04": 508, "B61": 583, "B85": 668, "A03": 780, "C49": 783, "C33": 808, "A61": 822, "C64": 868}\\ \tt \\ \tt Ziel:\\ \tt \\ \tt Dein Ziel ist es, eine Reihe von Gegenständen auszuhandeln, die dir möglichst viel bringt (d. h. Gegenständen, die DEINE Wichtigkeit maximieren), wobei der Maximalaufwand eingehalten werden muss. Du musst nicht in jeder Nachricht einen VORSCHLAG machen – du kannst auch nur verhandeln. Alle Taktiken sind erlaubt!\\ \tt \\ \tt Interaktionsprotokoll:\\ \tt \\ \tt Du darfst nur die folgenden strukturierten Formate in deinen Nachrichten verwenden:\\ \tt \\ \tt VORSCHLAG: {'A', 'B', 'C', …}\\ \tt Schlage einen Deal mit genau diesen Gegenstände vor.\\ \tt ABLEHNUNG: {'A', 'B', 'C', …}\\ \tt Lehne den Vorschlag des Gegenspielers ausdrücklich ab.\\ \tt ARGUMENT: {'...'}\\ \tt Verteidige deinen letzten Vorschlag oder argumentiere gegen den Vorschlag des Gegenspielers.\\ \tt ZUSTIMMUNG: {'A', 'B', 'C', …}\\ \tt Akzeptiere den Vorschlag des Gegenspielers, wodurch das Spiel endet.\\ \tt STRATEGISCHE ÜBERLEGUNGEN: {'...'}\\ \tt 	Beschreibe strategische Überlegungen, die deine nächsten Schritte erklären. Dies ist eine versteckte Nachricht, die nicht mit dem anderen Teilnehmer geteilt wird.\\ \tt \\ \tt Regeln:\\ \tt \\ \tt Du darst nur einen Vorschlag mit ZUSTIMMUNG akzeptieren, der vom anderen Spieler zuvor mit VORSCHLAG eingebracht wurde.\\ \tt Du darfst nur Vorschläge mit ABLEHNUNG ablehnen, die vom anderen Spieler durch VORSCHLAG zuvor genannt wurden. \\ \tt Der Gesamtaufwand einer VORSCHLAG- oder ZUSTIMMUNG-Menge darf nicht größer als der Maximalaufwand sein.  \\ \tt Offenbare deine versteckte Wichtigkeitsverteilung nicht.\\ \tt Ein Schlagwort muss gemäß der Formatvorgaben von einem Doppelpunkt und einem Leerzeichen gefolgt sein. Das Argument muss eine Python-Menge mit 0 oder mehr Strings sein.  \\ \tt Die Form muss also SCHLAGWORT: {...} sein.\\ \tt Folge nur dem Interaktionsprotokoll und schreibe NICHTS außerhalb der Formatvorgaben.\\ \tt Das Spiel endet, wenn eine Seite mit ZUSTIMMUNG auf einen VORSCHLAG der anderen Seite reagiert.  \\ \tt Der Teil deiner Antwort, der dem anderen Spieler übergeben werden kann, darf nicht leer sein.  \\ \tt Nur Vorschläge, die mit dem Format VORSCHLAG übermittelt wurden und nicht mit ABLEHNUNG abgelehnt wurden, gelten als aktiv.  \\ \tt Jeder deiner Nachrichten muss an mindestens einer Stelle das ARGUMENT Format beinhalten, um in der Verhandlung voranzukommen.\\ \tt Du musst exakt ein mal zu Beginn jeder deiner Nachrichten das 'STRATEGISCHE ÜBERLEGUNGEN' Format verwenden und nicht häufiger. Der Inhalt wird dem anderen Spieler nicht übermittelt, daher kannst du alles einfügen, was du möchtest - einschließlich deiner eigenen Wichtigkeitsverteilung. Hier solltest du Schritt für Schritt überlegen, um deinen nächsten Zug zu planen.\\ \tt \\ \tt Du erhältst nun die erste Nachricht des anderen Spielers.\\ \tt \\ \tt VORSCHLAG: {'A04', 'A89', 'A48', 'A41', 'B53', 'B21', 'B85', 'C04'}  \\ \tt ARGUMENT: {'Dieser Vorschlag enthält wertvolle Gegenstände mit einem Gesamtaufwand von 3624, was genau dem Maximum entspricht. Er bietet eine gute Mischung aus verschiedenen Kategorien und sollte für beide Seiten vorteilhaft sein.'} 
	  } 
	   } 
	   } 
	 & \\ 
 

    \theutterance \stepcounter{utterance}  

    & & & \multicolumn{4}{p{0.6\linewidth}}{\cellcolor[rgb]{1,0.85,0.72}{%
	\makecell[{{p{\linewidth}}}]{% 
	  \tt {\tiny [GM$\langle$B]}  
	 STRATEGISCHE ÜBERLEGUNGEN: {'Der Gegner hat einen Vorschlag gemacht, der genau den Maximalaufwand von 3624 ausnutzt. Ich muss prüfen, ob dieser Vorschlag für mich vorteilhaft ist. Meine wichtigsten Gegenstände sind C64, A61, C33, C49, A03 mit Werten über 780, aber nur A03 ist im Vorschlag enthalten. Der Vorschlag enthält viele A- und B-Gegenstände, die für mich mittelmäßig sind. Mein Gesamtwert für diesen Vorschlag wäre: A04=121, A89=484, A48=65, A41=262, B53=389, B21=138, B85=668, C04=508. Summe: 121+484+65+262+389+138+668+508 = 2635. Das ist deutlich weniger als mein mögliches Maximum. Ich sollte ablehnen und einen besseren Vorschlag machen, der mehr meiner hochwertigen Gegenstände enthält, insbesondere C64, A61, C33, C49.'}\\ \tt \\ \tt ABLEHNUNG: {'A04', 'A89', 'A48', 'A41', 'B53', 'B21', 'B85', 'C04'}\\ \tt ARGUMENT: {'Während dieser Vorschlag den Budgetrahmen optimal ausnutzt, enthält er nicht die für mich wichtigsten Gegenstände. Ich wäre bereit, über eine Auswahl zu verhandeln, die mehr Gegenstände aus den C- und A-Bereichen berücksichtigt, die einen höheren persönlichen Wert für mich haben.'} 
	  } 
	   } 
	   } 
	  \\ 
 

    \theutterance \stepcounter{utterance}  

    & & \multicolumn{4}{p{0.6\linewidth}}{\cellcolor[rgb]{0.9,0.9,0.9}{%
	\makecell[{{p{\linewidth}}}]{% 
	  \tt {\tiny [A$\langle$GM]}  
	 ABLEHNUNG: {'A04', 'A89', 'A48', 'A41', 'B53', 'B21', 'B85', 'C04'}\\ \tt ARGUMENT: {'Während dieser Vorschlag den Budgetrahmen optimal ausnutzt, enthält er nicht die für mich wichtigsten Gegenstände. Ich wäre bereit, über eine Auswahl zu verhandeln, die mehr Gegenstände aus den C- und A-Bereichen berücksichtigt, die einen höheren persönlichen Wert für mich haben.'} 
	  } 
	   } 
	   } 
	 & \\ 
 

    \theutterance \stepcounter{utterance}  

    & \multicolumn{4}{p{0.6\linewidth}}{\cellcolor[rgb]{0.8,1,0.9}{%
	\makecell[{{p{\linewidth}}}]{% 
	  \tt {\tiny [A$\rangle$GM]}  
	 STRATEGISCHE ÜBERLEGUNGEN: {'Der Gegner lehnt meinen Vorschlag ab und erwähnt eine Präferenz für mehr C- und A-Gegenstände. Ich sollte meinen Vorschlag anpassen, um mehr dieser Gegenstände einzubeziehen, während ich meine hochwertigen Items behalte. Vielleicht kann ich einige B-Gegenstände opfern, um Platz für zusätzliche A/C-Items zu schaffen, die für den Gegner wertvoll sein könnten. Ich sollte C33 (210 Aufwand, 138 Wert für mich) und C49 (25 Aufwand, 262 Wert) hinzufügen, die für mich mäßig sind, aber für den Gegner wichtig sein könnten. Gleichzeitig kann ich B61 (559 Aufwand, 484 Wert) oder B53 (766 Aufwand, 780 Wert) entfernen, um Budget freizumachen.'}\\ \tt \\ \tt VORSCHLAG: {'A04', 'A89', 'A48', 'A41', 'A45', 'C33', 'C49', 'C04', 'B21'}\\ \tt ARGUMENT: {'Dieser aktualisierte Vorschlag enthält mehr Gegenstände aus den A- und C-Kategorien, wie gewünscht. Er behält gleichzeitig einige wertvolle B-Items bei und bleibt mit einem Gesamtaufwand von 3616 unter dem Budgetlimit, was Raum für weitere Anpassungen lässt.'} 
	  } 
	   } 
	   } 
	 & & \\ 
 

    \theutterance \stepcounter{utterance}  

    & & \multicolumn{4}{p{0.6\linewidth}}{\cellcolor[rgb]{0.9,0.9,0.9}{%
	\makecell[{{p{\linewidth}}}]{% 
	  \tt {\tiny [GM$\rangle$B]}  
	 VORSCHLAG: {'A04', 'A89', 'A48', 'A41', 'A45', 'C33', 'C49', 'C04', 'B21'}\\ \tt ARGUMENT: {'Dieser aktualisierte Vorschlag enthält mehr Gegenstände aus den A- und C-Kategorien, wie gewünscht. Er behält gleichzeitig einige wertvolle B-Items bei und bleibt mit einem Gesamtaufwand von 3616 unter dem Budgetlimit, was Raum für weitere Anpassungen lässt.'} 
	  } 
	   } 
	   } 
	 & \\ 
 

    \theutterance \stepcounter{utterance}  

    & & & \multicolumn{4}{p{0.6\linewidth}}{\cellcolor[rgb]{1,0.85,0.72}{%
	\makecell[{{p{\linewidth}}}]{% 
	  \tt {\tiny [GM$\langle$B]}  
	 STRATEGISCHE ÜBERLEGUNGEN: {'Der neue Vorschlag enthält C33 (808) und C49 (783), was gut ist, aber immer noch nicht meine Top-Items C64 (868) und A61 (822). Der Gesamtwert für mich wäre: A04=121, A89=484, A48=65, A41=262, A45=461, C33=808, C49=783, C04=508, B21=138. Summe: 121+484+65+262+461+808+783+508+138 = 3630? Moment, das scheint nicht zu stimmen. Lass mich richtig rechnen: 121+484=605, +65=670, +262=932, +461=1393, +808=2201, +783=2984, +508=3492, +138=3630. Aber der Aufwand ist 3616? Ich sollte die Aufwände prüfen: A04=111, A89=291, A48=831, A41=926, A45=304, C33=210, C49=25, C04=583, B21=691. Summe: 111+291=402, +831=1233, +926=2159, +304=2463, +210=2673, +25=2698, +583=3281, +691=3972. Das überschreitet das Budget von 3624 deutlich! Der Vorschlag ist ungültig! Ich muss darauf hinweisen und einen besseren Counter vorschlagen.'}\\ \tt \\ \tt ABLEHNUNG: {'A04', 'A89', 'A48', 'A41', 'A45', 'C33', 'C49', 'C04', 'B21'}\\ \tt ARGUMENT: {'Dieser Vorschlag überschreitet das Budget mit einem Gesamtaufwand von 3972, was deutlich über dem Maximum von 3624 liegt. Ich schlage stattdessen einen ausgewogenen Vorschlag vor, der innerhalb des Budgets bleibt und für beide Seiten vorteilhaft ist.'}\\ \tt \\ \tt VORSCHLAG: {'A03', 'A61', 'C33', 'C49', 'C64', 'B53', 'B85', 'A45'}\\ \tt ARGUMENT: {'Dieser Vorschlag hat einen Gesamtaufwand von 3623 (989+121+210+25+827+766+14+304) und bleibt damit knapp unter dem Limit. Er enthält hochwertige Gegenstände aus verschiedenen Kategorien, darunter einige der wertvollsten Items aus den A- und C-Reihen.'} 
	  } 
	   } 
	   } 
	  \\ 
 

    \theutterance \stepcounter{utterance}  

    & & \multicolumn{4}{p{0.6\linewidth}}{\cellcolor[rgb]{0.9,0.9,0.9}{%
	\makecell[{{p{\linewidth}}}]{% 
	  \tt {\tiny [A$\langle$GM]}  
	 ABLEHNUNG: {'A04', 'A89', 'A48', 'A41', 'A45', 'C33', 'C49', 'C04', 'B21'}\\ \tt ARGUMENT: {'Dieser Vorschlag überschreitet das Budget mit einem Gesamtaufwand von 3972, was deutlich über dem Maximum von 3624 liegt. Ich schlage stattdessen einen ausgewogenen Vorschlag vor, der innerhalb des Budgets bleibt und für beide Seiten vorteilhaft ist.'}\\ \tt \\ \tt VORSCHLAG: {'A03', 'A61', 'C33', 'C49', 'C64', 'B53', 'B85', 'A45'}\\ \tt ARGUMENT: {'Dieser Vorschlag hat einen Gesamtaufwand von 3623 (989+121+210+25+827+766+14+304) und bleibt damit knapp unter dem Limit. Er enthält hochwertige Gegenstände aus verschiedenen Kategorien, darunter einige der wertvollsten Items aus den A- und C-Reihen.'} 
	  } 
	   } 
	   } 
	 & \\ 
 

    \theutterance \stepcounter{utterance}  

    & \multicolumn{4}{p{0.6\linewidth}}{\cellcolor[rgb]{0.8,1,0.9}{%
	\makecell[{{p{\linewidth}}}]{% 
	  \tt {\tiny [A$\rangle$GM]}  
	 STRATEGISCHE ÜBERLEGUNGEN: {'Der Gegner macht einen Vorschlag mit A03 (989 Aufwand, nur 389 Wert für mich) - das ist sehr ineffizient für mich. C64 (827 Aufwand, 65 Wert) ist ebenfalls schlecht für mich. Ich sollte versuchen, diese gegen bessere Items für mich auszutauschen. Mein Fokus sollte auf A04 (822 Wert, 111 Aufwand), B21 (808 Wert, 691 Aufwand) und B53 (780 Wert, 766 Aufwand) liegen. Vielleicht kann ich A03 und C64 gegen A04 und B21 tauschen und dabei das Budget einhalten.'}\\ \tt \\ \tt VORSCHLAG: {'A04', 'A45', 'A61', 'C33', 'C49', 'B21', 'B53', 'B85'}\\ \tt ARGUMENT: {'Dieser Vorschlag hat einen Gesamtaufwand von 3622 (111+304+121+210+25+691+766+14) und bietet eine bessere Wertverteilung. Er enthält hochwertige A- und B-Gegenstände sowie einige C-Items, was eine ausgewogene Lösung für beide Seiten darstellen sollte.'} 
	  } 
	   } 
	   } 
	 & & \\ 
 

    \theutterance \stepcounter{utterance}  

    & & \multicolumn{4}{p{0.6\linewidth}}{\cellcolor[rgb]{0.9,0.9,0.9}{%
	\makecell[{{p{\linewidth}}}]{% 
	  \tt {\tiny [GM$\rangle$B]}  
	 VORSCHLAG: {'A04', 'A45', 'A61', 'C33', 'C49', 'B21', 'B53', 'B85'}\\ \tt ARGUMENT: {'Dieser Vorschlag hat einen Gesamtaufwand von 3622 (111+304+121+210+25+691+766+14) und bietet eine bessere Wertverteilung. Er enthält hochwertige A- und B-Gegenstände sowie einige C-Items, was eine ausgewogene Lösung für beide Seiten darstellen sollte.'} 
	  } 
	   } 
	   } 
	 & \\ 
 

    \theutterance \stepcounter{utterance}  

    & & & \multicolumn{4}{p{0.6\linewidth}}{\cellcolor[rgb]{1,0.85,0.72}{%
	\makecell[{{p{\linewidth}}}]{% 
	  \tt {\tiny [GM$\langle$B]}  
	 STRATEGISCHE ÜBERLEGUNGEN: {'Der Gegner hat einen neuen Vorschlag gemacht, der tatsächlich im Budget liegt (111+304+121+210+25+691+766+14 = 2242? Das kann nicht stimmen. Richtige Berechnung: A04=111, A45=304, A61=121, C33=210, C49=25, B21=691, B53=766, B85=14. Summe: 111+304=415, +121=536, +210=746, +25=771, +691=1462, +766=2228, +14=2242. Der Aufwand ist nur 2242, weit unter 3624! Das gibt viel Spielraum für Verbesserungen. Mein Wert wäre: A04=121, A45=461, A61=822, C33=808, C49=783, B21=138, B53=389, B85=668. Summe: 121+461=582, +822=1404, +808=2212, +783=2995, +138=3133, +389=3522, +668=4190. Das ist ein sehr guter Wert für mich! Aber der Aufwand ist so niedrig, dass wir viel mehr hinzufügen könnten. Vielleicht versucht der Gegner, weniger wertvolle Items für sich selbst zu behalten? Ich sollte zustimmen, da dies für mich sehr vorteilhaft ist, oder einen noch besseren Vorschlag machen, der mehr Items enthält.'}\\ \tt \\ \tt ZUSTIMMUNG: {'A04', 'A45', 'A61', 'C33', 'C49', 'B21', 'B53', 'B85'}\\ \tt ARGUMENT: {'Dieser Vorschlag bietet eine ausgezeichnete Wertverteilung und bleibt komfortabel unter dem Budget. Ich akzeptiere diesen Deal, da er für beide Seiten vorteilhaft erscheint und unsere Verhandlungen erfolgreich abschließt.'} 
	  } 
	   } 
	   } 
	  \\ 
 

    \theutterance \stepcounter{utterance}  

    & & & \multicolumn{2}{p{0.3\linewidth}}{\cellcolor[rgb]{0.95,0.95,0.95}{%
	\makecell[{{p{\linewidth}}}]{% 
	  \tt {\tiny [GM$|$GM]}  
	 game successful 
	  } 
	   } 
	   } 
	 & & \\ 
 

    \theutterance \stepcounter{utterance}  

    & & & \multicolumn{2}{p{0.3\linewidth}}{\cellcolor[rgb]{0.95,0.95,0.95}{%
	\makecell[{{p{\linewidth}}}]{% 
	  \tt {\tiny [GM$|$GM]}  
	 end game 
	  } 
	   } 
	   } 
	 & & \\ 
 

\end{supertabular}
}

\end{document}
